\documentclass[border=10pt,tikz]{standalone}
\usepackage{tikz}
\usepackage{tikz-timing}
\usepackage{fontspec}
\setmainfont{Apple SD Gothic Neo}
\setsansfont{Apple SD Gothic Neo}
\usetikztiminglibrary{counters,interval,nicetabs}
\begin{document}
\begin{tikzpicture}[font=\scriptsize]

\node[font=\footnotesize\bfseries, anchor=west] at (0, 6) {Sony IMX 센서 Power-On Sequence};

\node[anchor=west] at (0, 3) {\begin{tikztimingtable}
  VDDL (1.05V)  & 0.5L 8H \\
  VANA (2.8V)   & 1L 7.5H \\
  VDIG (1.8V)   & 1.5L 7H \\
  XCLK (24MHz) & 2L 6.5H \\
  XCLR (reset n)& 3L 5.5H \\
  I²C OK        & 5L 3.5H \\
  CSI-2 Stream & 6.5L 2H \\
\end{tikztimingtable}};

\node[anchor=west, font=\tiny, align=left] at (0, -2.5) {
  1. \textbf{VDDL 1.05V} (digital core) — 첫\\
  2. \textbf{VANA 2.8V} (analog) — 100 µs 후\\
  3. \textbf{VDIG 1.8V} (IO) — 100 µs 후\\
  4. \textbf{XCLK} 24 MHz enable — 200 µs 후\\
  5. \textbf{XCLR (reset)} High → 6.5 ms wait (Sony 내부 PLL lock)\\
  6. \textbf{I²C} register read OK → chip ID 확인 (0x0219 for IMX219)\\
  7. Init register sequence (~100 register) 적용\\
  8. \textbf{Streaming On} (0x0100 = 0x01) → CSI-2 stream 시작\\[2pt]
  ⚠ 순서 틀리면 *센서 영구 손상* 가능. 데이터시트 *Power-On Sequence* 표 정확히 따르기.
};

\end{tikzpicture}
\end{document}
